% GENERATED FILE. Edit canonical lessons, then run npm run generate:books.
% canonical-source-hash: fnv1a64:f44dc60e7e53ab54
% canonical-lessons: TE-C378-kotlata, TE-C378-kalpulu, TE-C378-gayaparacu, TE-C378-penuguladu, TE-C378-gumigudu

\chapter{A Brawl}
\label{ch:te-a-brawl}
\hlchaptermodality{378}

\begin{chapteropening}
\textbf{By the end of this chapter:} \emph{I can say a brawl, gunfire, to injure, to grapple and to gather round.}

Complete the last lesson of chapter 378: I can say a brawl, gunfire, to injure, to grapple and to gather round.
\end{chapteropening}

\section[\texorpdfstring{\te{కొట్లాట} (koṭlāṭa)}{కొట్లాట (koṭlāṭa)}]{\te{కొట్లాట} (koṭlāṭa) — a brawl; a scuffle}

\label{lesson:TE-C378-kotlata}



\begin{quote}
Before the new one: say the Telugu for a patrol, then the Telugu for an action.
\end{quote}

\subsection*{You'll want to know: \te{కొట్లాట}}
\textbf{\te{కొట్లాట}} — \emph{koṭlāṭa} — \textquotedblleft{}a brawl; a scuffle\textquotedblright{}.

\begin{quote}
\te{బజారులో} \te{ఇద్దరి} \te{మధ్య} \te{కొట్లాట} \te{జరిగింది}.
\end{quote}

\begin{quote}
\emph{bajārulō iddari madhya koṭlāṭa jarigindi.}
\end{quote}

\begin{quote}
There was a brawl between two men in the market.
\end{quote}

Plural: \te{కొట్లాటలు}, \emph{koṭlāṭalu}.

From \te{కొట్టు}, to hit: a fight with blows.

\begin{cousinweb}
\textbf{\te{కొట్లాట}} (\emph{koṭlāṭa}) is native Dravidian, built on \emph{koṭṭu}, \textquotedblleft{}to beat, strike\textquotedblright{} (DEDR 2063).
\end{cousinweb}

\begin{grammarlens}[title={What you've built}]
One more word for telling what happened.
\end{grammarlens}

\subsection*{Your turn}
\begin{itemize}
\raggedright
  \item \emph{Say it:} \emph{koṭlāṭa}
  \item \emph{Say it:} \emph{koṭlāṭa}, once more
  \item \emph{Say it:} \emph{bajārulō iddari madhya koṭlāṭa jarigindi.} — There was a brawl between two men in the market.
  \item \emph{Recall:} say the Telugu for a patrol, then the Telugu for an action, then say \emph{koṭlāṭa} again
\end{itemize}

\begin{tcolorbox}[breakable,colback=teal!4,colframe=teal!35!black,title={Before you move on}]
What does \te{కొట్లాట} mean? (\textquotedblleft{}A brawl; a scuffle\textquotedblright{}.) Say it once more.
\end{tcolorbox}
\section[\texorpdfstring{\te{కాల్పులు} (kālpulu)}{కాల్పులు (kālpulu)}]{\te{కాల్పులు} (kālpulu) — gunfire; shooting}

\label{lesson:TE-C378-kalpulu}



\begin{quote}
Before the new one: say the Telugu for an action, then the Telugu for a brawl.
\end{quote}

\subsection*{You'll want to know: \te{కాల్పులు}}
\textbf{\te{కాల్పులు}} — \emph{kālpulu} — \textquotedblleft{}gunfire; shooting\textquotedblright{}.

\begin{quote}
\te{రాత్రి} \te{అడవిలో} \te{కాల్పులు} \te{వినిపించాయి}.
\end{quote}

\begin{quote}
\emph{rātri aḍavilō kālpulu vinipiñcāyi.}
\end{quote}

\begin{quote}
Gunfire was heard in the forest at night.
\end{quote}

A plural noun from \te{కాల్చు}, to burn or to fire a gun.

\begin{cousinweb}
\textbf{\te{కాల్పులు}} (\emph{kālpulu}) is native Dravidian, the plural noun of \emph{kālcu}, \textquotedblleft{}to burn, set fire to\textquotedblright{} (DEDR 1500), the verb also used of firing a gun.
\end{cousinweb}

\begin{grammarlens}[title={What you've built}]
One more word for telling what happened.
\end{grammarlens}

\subsection*{Your turn}
\begin{itemize}
\raggedright
  \item \emph{Say it:} \emph{kālpulu}
  \item \emph{Say it:} \emph{kālpulu}, once more
  \item \emph{Say it:} \emph{rātri aḍavilō kālpulu vinipiñcāyi.} — Gunfire was heard in the forest at night.
  \item \emph{Recall:} say the Telugu for an action, then the Telugu for a brawl, then say \emph{kālpulu} again
\end{itemize}

\begin{tcolorbox}[breakable,colback=teal!4,colframe=teal!35!black,title={Before you move on}]
What does \te{కాల్పులు} mean? (\textquotedblleft{}Gunfire; shooting\textquotedblright{}.) Say it once more.
\end{tcolorbox}
\section[\texorpdfstring{\te{గాయపరచు} (gāyaparacu)}{గాయపరచు (gāyaparacu)}]{\te{గాయపరచు} (gāyaparacu) — to injure; to wound}

\label{lesson:TE-C378-gayaparacu}



\begin{quote}
Before the new one: say the Telugu for a brawl, then the Telugu for gunfire.
\end{quote}

\subsection*{You'll want to know: \te{గాయపరచు}}
\textbf{\te{గాయపరచు}} — \emph{gāyaparacu} — \textquotedblleft{}to injure; to wound\textquotedblright{}.

\begin{quote}
\te{దొంగ} \te{కాపలాదారును} \te{కత్తితో} \te{గాయపరిచాడు}.
\end{quote}

\begin{quote}
\emph{doṅga kāpalādārunu kattitō gāyaparicāḍu.}
\end{quote}

\begin{quote}
The thief injured the watchman with a knife.
\end{quote}

Past (I …): \te{గాయపరిచాను}, \emph{gāyaparicānu}.

You know \te{గాయం}, a wound; the ending -\te{పరచు} turns it into something done to another.

\begin{cousinweb}
\textbf{\te{గాయపరచు}} (\emph{gāyaparacu}) is built on \te{గాయం}, \textquotedblleft{}a wound\textquotedblright{}, which goes back to Sanskrit \emph{ghāta}, \textquotedblleft{}a blow, a bruise\textquotedblright{} (CDIAL 4460), with \emph{paracu}, which makes a verb of a noun.
\end{cousinweb}

\begin{grammarlens}[title={What you've built}]
One more word for telling what happened.
\end{grammarlens}

\subsection*{Your turn}
\begin{itemize}
\raggedright
  \item \emph{Say it:} \emph{gāyaparacu}
  \item \emph{Say it:} \emph{gāyaparacu}, once more
  \item \emph{Say it:} \emph{doṅga kāpalādārunu kattitō gāyaparicāḍu.} — The thief injured the watchman with a knife.
  \item \emph{Recall:} say the Telugu for a brawl, then the Telugu for gunfire, then say \emph{gāyaparacu} again
\end{itemize}

\begin{tcolorbox}[breakable,colback=teal!4,colframe=teal!35!black,title={Before you move on}]
What does \te{గాయపరచు} mean? (\textquotedblleft{}To injure; to wound\textquotedblright{}.) Say it once more.
\end{tcolorbox}
\section[\texorpdfstring{\te{పెనుగులాడు} (penugulāḍu)}{పెనుగులాడు (penugulāḍu)}]{\te{పెనుగులాడు} (penugulāḍu) — to grapple; to wrestle free}

\label{lesson:TE-C378-penuguladu}



\begin{quote}
Before the new one: say the Telugu for gunfire, then the Telugu for to injure.
\end{quote}

\subsection*{You'll want to know: \te{పెనుగులాడు}}
\textbf{\te{పెనుగులాడు}} — \emph{penugulāḍu} — \textquotedblleft{}to grapple; to wrestle free\textquotedblright{}.

\begin{quote}
\te{ఆమె} \te{దొంగతో} \te{పెనుగులాడింది}.
\end{quote}

\begin{quote}
\emph{āme doṅgatō penugulāḍindi.}
\end{quote}

\begin{quote}
She grappled with the thief.
\end{quote}

Past (I …): \te{పెనుగులాడాను}, \emph{penugulāḍānu}.

With whom takes -\te{తో}.

\begin{cousinweb}
\textbf{\te{పెనుగులాడు}} (\emph{penugulāḍu}) is native Dravidian: DEDR 4160 lists Telugu \emph{penãgu}, \textquotedblleft{}to be twisted, to struggle, fight\textquotedblright{}, here with \emph{āḍu}.
\end{cousinweb}

\begin{grammarlens}[title={What you've built}]
One more word for telling what happened.
\end{grammarlens}

\subsection*{Your turn}
\begin{itemize}
\raggedright
  \item \emph{Say it:} \emph{penugulāḍu}
  \item \emph{Say it:} \emph{penugulāḍu}, once more
  \item \emph{Say it:} \emph{āme doṅgatō penugulāḍindi.} — She grappled with the thief.
  \item \emph{Recall:} say the Telugu for gunfire, then the Telugu for to injure, then say \emph{penugulāḍu} again
\end{itemize}

\begin{tcolorbox}[breakable,colback=teal!4,colframe=teal!35!black,title={Before you move on}]
What does \te{పెనుగులాడు} mean? (\textquotedblleft{}To grapple; to wrestle free\textquotedblright{}.) Say it once more.
\end{tcolorbox}
\section[\texorpdfstring{\te{గుమిగూడు} (gumigūḍu)}{గుమిగూడు (gumigūḍu)}]{\te{గుమిగూడు} (gumigūḍu) — to gather round}

\label{lesson:TE-C378-gumigudu}



\begin{quote}
Before the new one: say the Telugu for to injure, then the Telugu for to grapple.
\end{quote}

\subsection*{You'll want to know: \te{గుమిగూడు}}
\textbf{\te{గుమిగూడు}} — \emph{gumigūḍu} — \textquotedblleft{}to gather round\textquotedblright{}.

\begin{quote}
\te{ప్రమాదం} \te{జరిగిన} \te{చోట} \te{జనం} \te{గుమిగూడారు}.
\end{quote}

\begin{quote}
\emph{pramādaṁ jarigina cōṭa janaṁ gumigūḍāru.}
\end{quote}

\begin{quote}
People gathered where the accident happened.
\end{quote}

Past (I …): \te{గుమిగూడాను}, \emph{gumigūḍānu}.

\te{గుమి}, a crowd, and \te{కూడు}, to come together.

\begin{cousinweb}
\textbf{\te{గుమిగూడు}} (\emph{gumigūḍu}) is a native compound: \emph{gumi}, \textquotedblleft{}a crowd\textquotedblright{} (DEDR 1741), and \emph{kūḍu}, \textquotedblleft{}to meet, gather\textquotedblright{} (DEDR 1882).
\end{cousinweb}

\begin{grammarlens}[title={What you've built}]
One more word for telling what happened.
\end{grammarlens}

\subsection*{Your turn}
\begin{itemize}
\raggedright
  \item \emph{Say it:} \emph{gumigūḍu}
  \item \emph{Say it:} \emph{gumigūḍu}, once more
  \item \emph{Say it:} \emph{pramādaṁ jarigina cōṭa janaṁ gumigūḍāru.} — People gathered where the accident happened.
  \item \emph{Recall:} say the Telugu for to injure, then the Telugu for to grapple, then say \emph{gumigūḍu} again
\end{itemize}

\begin{tcolorbox}[breakable,colback=teal!4,colframe=teal!35!black,title={Before you move on}]
What does \te{గుమిగూడు} mean? (\textquotedblleft{}To gather round\textquotedblright{}.) Say it once more.
\end{tcolorbox}
